\documentclass[12pt,a4paper]{article}

\usepackage[spanish]{babel}
\usepackage[utf8]{inputenc}
\usepackage{amsmath, amssymb, amsthm}
\usepackage{graphicx}
\usepackage{booktabs}
\usepackage{geometry}
\usepackage{physics}
\usepackage{multicol}
\usepackage{hyperref}
\usepackage{float}

\geometry{left=3cm, right=3cm, top=3cm, bottom=3cm}

\title{Integración Numérica}
\author{Métodos Numéricos II}
\date{}

\begin{document}
\maketitle


\section*{Calendario de clases – Abril y Mayo}

\begin{center}
\begin{tabular}{ccc}
\toprule
Semana & Martes & Viernes \\
\midrule
Semana 1 & 14 Abril & 17 Abril \\
Semana 2 & 21 Abril & 24 Abril \\
Semana 3 & 28 Abril & 1 Mayo \\
Semana 4 & 5 Mayo & 8 Mayo \\
Semana 5 & 12 Mayo & 15 Mayo \\
\bottomrule
\end{tabular}
\end{center}


\section{¿Por qué necesitamos las integrales numéricas?}

En muchos problemas reales de ingeniería, física y matemáticas aplicadas, las integrales definidas no pueden calcularse de forma exacta por varias razones:

\begin{itemize}
\item La función no tiene primitiva elemental.
\item La función está dada por datos experimentales.
\item La función es muy complicada para integrar analíticamente.
\item Solo conocemos valores de la función en puntos discretos.
\end{itemize}

Por ello, se utilizan métodos de \textbf{integración numérica}, que aproximan el valor de la integral:

\[
\int_a^b f(x)\,dx
\]

mediante sumas ponderadas de valores de la función.\\

{\it 
La idea general de las reglas de integración numérica es aproximar la función por un polinomio interpolador y luego integrar ese polinomio.}

\section{Métodos de integración numérica}

El objetivo de la integración numérica es aproximar integrales definidas de la forma:

\[
\int_a^b f(x)\,dx
\]

cuando la integral no puede calcularse de forma exacta o cuando la función solo se conoce en puntos discretos.

Los métodos de integración numérica más importantes se clasifican en dos grandes grupos:

\begin{itemize}
\item Métodos de Newton-Cotes
\item Cuadratura de Gauss
\end{itemize}

\subsection{Métodos de Newton-Cotes}

Los métodos de Newton-Cotes consisten en aproximar la función $f(x)$ mediante un 
polinomio interpolador construido usando puntos equiespaciados, y luego integrar
dicho polinomio.

Es decir:

\[
\int_a^b f(x)\,dx \approx \int_a^b p_n(x)\,dx
\]

donde $p_n(x)$ es el polinomio interpolador que pasa por puntos equiespaciados.

Los métodos más conocidos de Newton-Cotes son:

\begin{itemize}
\item Regla del rectángulo
\item Regla del trapecio
\item Regla de Simpson 1/3
\item Regla de Simpson 3/8
\end{itemize}

En estos métodos los puntos están igualmente espaciados:

\[
x_i = a + ih
\]

\subsection{Cuadratura de Gauss}

La cuadratura de Gauss también aproxima la integral mediante una suma ponderada de valores
de la función:

\[
\int_a^b f(x)\,dx \approx \sum_{i=1}^{n} w_i f(x_i)
\]

Sin embargo, a diferencia de Newton-Cotes, los puntos $x_i$ no son equiespaciados,
sino que se eligen de forma óptima para maximizar la precisión del método.

La cuadratura de Gauss con $n$ puntos es exacta para polinomios de grado $2n-1$,
mientras que los métodos de Newton-Cotes con $n$ puntos son exactos para polinomios
de grado $n$ o menor.

\subsection{Diferencia entre Newton-Cotes y Gauss}

\begin{center}
\begin{tabular}{lcc}
\toprule
Método & Puntos & Precisión \\
\midrule
Newton-Cotes & Equiespaciados & Menor \\
Gauss & No equiespaciados & Mayor \\
\bottomrule
\end{tabular}
\end{center}

En resumen:

\begin{itemize}
\item Newton-Cotes: puntos equiespaciados.
\item Gauss: puntos óptimos.
\item Gauss suele ser más preciso con menos puntos.
\end{itemize}


\subsection{Precisión de un método numérico}

La precisión de un método numérico mide qué tan cerca está la aproximación numérica
del valor exacto.

Se define el error como:

\[
\text{Error} = \left| I - I_h \right|
\]

donde:
\begin{itemize}
\item $I$ es el valor exacto de la integral
\item $I_h$ es la aproximación numérica
\item $h$ es el tamaño de paso
\end{itemize}

Decimos que un método numérico tiene orden de precisión $p$ si el error se comporta como:

\[
E(h) = C h^p
\]

donde:
\begin{itemize}
\item $C$ es una constante
\item $h$ es el tamaño de paso
\item $p$ es el orden del método
\end{itemize}

Esto significa que cuando el tamaño de paso disminuye, el error disminuye según la potencia $h^p$.

\subsection{Interpretación}

Si un método es de orden $p$, entonces:

\begin{itemize}
\item Si $h$ se reduce a la mitad, el error se reduce aproximadamente por un factor de $2^p$.
\end{itemize}

\begin{center}
\begin{tabular}{ccc}
\toprule
Método & Orden & Reducción del error si $h/2$ \\
\midrule
Trapecio & $p=2$ & Error se divide entre 4 \\
Simpson & $p=4$ & Error se divide entre 16 \\
Gauss (2 puntos) & $p=4$ & Error se divide entre 16 \\
\bottomrule
\end{tabular}
\end{center}


Key: Un método es más preciso si su orden $p$ es mayor, ya que el error disminuye más rápido
cuando el tamaño de paso $h$ disminuye.

\newpage
\section{Principales métodos de integración numérica}

\begin{table}[h!]
\centering
\begin{tabular}{lccc}
\toprule
Método & Tipo de aproximación & Orden del error & Observaciones \\
\midrule
Rectángulo & Constante & $O(h)$ & Muy simple, poco preciso \\
Trapecio & Lineal & $O(h^2)$ & Fácil y estable \\
Simpson 1/3 & Cuadrática & $O(h^4)$ & Muy preciso \\
Simpson 3/8 & Cúbica & $O(h^4)$ & Variante de Simpson \\
Newton-Cotes & Polinomios & Depende del grado & Generalización \\
Cuadratura Gauss & Polinomios óptimos & Muy alto & Muy eficiente \\
\bottomrule
\end{tabular}
\caption{Métodos principales de integración numérica}
\end{table}

\section{Semejanzas y diferencias}

\subsection{Semejanzas}
\begin{itemize}
\item Todos aproximan la integral mediante sumas ponderadas.
\item Todos se basan en interpolación polinómica.
\item El error depende del tamaño de paso $h$.
\end{itemize}

\subsection{Diferencias}
\begin{itemize}
\item El grado del polinomio interpolador.
\item La precisión (orden del error).
\item El número de puntos utilizados.
\item La complejidad computacional.
\end{itemize}

\section{Regla del Trapecio}

\subsection{Deducción}

Aproximamos la función por un polinomio de grado 1 (recta) que pasa por los puntos:

\[
(x_0, f(x_0)), \quad (x_1, f(x_1))
\]

El polinomio interpolador es:

\[
p_1(x) = f(x_0)\frac{x - x_1}{x_0 - x_1} + f(x_1)\frac{x - x_0}{x_1 - x_0}
\]

Integramos el polinomio:

\[
\int_{x_0}^{x_1} f(x)\,dx \approx \int_{x_0}^{x_1} p_1(x)\,dx
\]

Después de integrar:

\[
\int_{x_0}^{x_1} f(x)\,dx \approx \frac{h}{2} \left[ f(x_0) + f(x_1) \right]
\]

donde $h = x_1 - x_0$.


\subsection{¿De dónde sale el polinomio $p_1(x)$?}

La regla del trapecio se basa en aproximar la función $f(x)$ por un 
\textbf{polinomio interpolador de grado 1} (una recta) que pasa por dos puntos.

Sea la integral:

\[
\int_a^b f(x)\,dx
\]

Tomamos dos puntos:

\[
x_0 = a, \quad x_1 = b
\]

Interpolamos la función mediante el polinomio de Lagrange de grado 1:

\[
p_1(x) = f(x_0)L_0(x) + f(x_1)L_1(x)
\]

donde las bases de Lagrange son:

\[
L_0(x) = \frac{x - x_1}{x_0 - x_1}, 
\quad
L_1(x) = \frac{x - x_0}{x_1 - x_0}
\]

Por lo tanto:

\[
p_1(x) = f(x_0)\frac{x - x_1}{x_0 - x_1}
+ f(x_1)\frac{x - x_0}{x_1 - x_0}
\]

Este polinomio es la recta que pasa por los puntos $(x_0, f(x_0))$ y $(x_1, f(x_1))$.

La idea del método del trapecio es aproximar:

\[
\int_a^b f(x)\,dx \approx \int_a^b p_1(x)\,dx
\]

Ahora integramos el polinomio:

\[
\int_a^b p_1(x)\,dx
=
\int_a^b
\left[
f(x_0)\frac{x - x_1}{x_0 - x_1}
+ f(x_1)\frac{x - x_0}{x_1 - x_0}
\right] dx
\]

Sacando constantes:

\[
=
f(x_0)\int_a^b \frac{x - x_1}{x_0 - x_1} dx
+
f(x_1)\int_a^b \frac{x - x_0}{x_1 - x_0} dx
\]

Resolviendo las integrales y simplificando, se obtiene:

\[
\int_a^b f(x)\,dx \approx \frac{b-a}{2}\left[f(a) + f(b)\right]
\]

Esta es la \textbf{regla del trapecio}.


\subsection{Error del Trapecio}

El error viene dado por:

\[
E_T = -\frac{(b-a)^3}{12} f''(\xi)
\]

para algún $\xi \in (a,b)$.

Orden del método:

\[
E_T = O(h^2)
\]


\subsection{Deducción del error de la regla del trapecio}

La regla del trapecio aproxima la integral de una función $f(x)$ en el intervalo $[a,b]$ 
mediante el polinomio interpolador de grado 1, $p_1(x)$:

\[
\int_a^b f(x)\,dx \approx \int_a^b p_1(x)\,dx
\]

Por lo tanto, el error viene dado por:

\[
E_T = \int_a^b f(x)\,dx - \int_a^b p_1(x)\,dx
= \int_a^b \left[f(x) - p_1(x)\right] dx
\]

\subsubsection*{Error del polinomio interpolador}

Sabemos por la teoría de interpolación de Lagrange que el error al interpolar con un 
polinomio de grado 1 es:

\[
f(x) - p_1(x) = \frac{f''(\xi_x)}{2!}(x-a)(x-b)
\]

para algún $\xi_x \in (a,b)$.

Sustituyendo en la expresión del error:

\[
E_T = \int_a^b \frac{f''(\xi_x)}{2}(x-a)(x-b)\,dx
\]

Si $f''$ es continua, podemos aplicar el teorema del valor medio para integrales, 
y obtenemos:

\[
E_T = \frac{f''(\xi)}{2} \int_a^b (x-a)(x-b)\,dx
\]

para algún $\xi \in (a,b)$.

\subsubsection*{Cálculo de la integral}

Calculamos:

\[
\int_a^b (x-a)(x-b)\,dx
\]

Primero expandimos:

\[
(x-a)(x-b) = x^2 - (a+b)x + ab
\]

Entonces:

\[
\int_a^b (x^2 - (a+b)x + ab)\,dx
\]

\[
= \left[ \frac{x^3}{3} - \frac{(a+b)x^2}{2} + abx \right]_a^b
\]

Evaluando en los límites y simplificando, se obtiene:

\[
\int_a^b (x-a)(x-b)\,dx = -\frac{(b-a)^3}{6}
\]

\subsubsection*{Error final}

Sustituyendo en la expresión del error:

\[
E_T = \frac{f''(\xi)}{2} \left(-\frac{(b-a)^3}{6}\right)
\]

Por lo tanto, el error de la regla del trapecio es:

\[
E_T = -\frac{(b-a)^3}{12} f''(\xi), \quad \xi \in (a,b)
\]

\subsubsection*{Orden del método}

Si definimos $h = b-a$, entonces:

\[
E_T = -\frac{h^3}{12} f''(\xi)
\]

Por lo tanto, el método del trapecio es de orden:

\[
E_T = O(h^2)
\]

El error depende de la segunda derivada de la función, es decir, de la curvatura de la función.

\subsection{Ejemplo}

Calcular:

\[
\int_0^1 x^2 dx
\]

Sabemos que el valor exacto es:

\[
\int_0^1 x^2 dx = \frac{1}{3}
\]

Aplicando trapecio:

\[
T = \frac{h}{2} [f(0) + f(1)]
\]

\[
T = \frac{1}{2} [0 + 1] = \frac{1}{2}
\]

Error:

\[
Error = \left| \frac{1}{3} - \frac{1}{2} \right| = 0.1667
\]




\newpage

\section{Regla de Simpson}

Simpson aproxima la función por un polinomio de grado 2.

\subsection{Simpson 1/3}

\[
\int_a^b f(x)\,dx \approx \frac{h}{3} \left[f(x_0) + 4f(x_1) + f(x_2)\right]
\]

Error:

\[
E_S = -\frac{(b-a)^5}{180} f^{(4)}(\xi)
\]

Orden:

\[
O(h^4)
\]
%%%%%%%%%%%%%%%%%%%%%%%%%%%%%%



\subsection{Deducción de la fórmula de Simpson 1/3}

La regla de Simpson se obtiene aproximando la función $f(x)$ mediante un 
polinomio interpolador de grado 2 (parábola).

Queremos aproximar:

\[
\int_a^b f(x)\,dx
\]

Dividimos el intervalo en dos subintervalos iguales:

\[
h = \frac{b-a}{2}
\]

Definimos los puntos:

\[
x_0 = a, \quad x_1 = a+h, \quad x_2 = b
\]

Interpolamos $f(x)$ mediante el polinomio de Lagrange de grado 2:

\[
p_2(x) = f(x_0)L_0(x) + f(x_1)L_1(x) + f(x_2)L_2(x)
\]

donde:

\[
L_0(x) = \frac{(x-x_1)(x-x_2)}{(x_0-x_1)(x_0-x_2)}
\]

\[
L_1(x) = \frac{(x-x_0)(x-x_2)}{(x_1-x_0)(x_1-x_2)}
\]

\[
L_2(x) = \frac{(x-x_0)(x-x_1)}{(x_2-x_0)(x_2-x_1)}
\]

La regla de Simpson consiste en aproximar:

\[
\int_a^b f(x)\,dx \approx \int_a^b p_2(x)\,dx
\]

Al integrar el polinomio interpolador y simplificar (cálculo algebraico), se obtiene:

\[
\int_a^b f(x)\,dx \approx \frac{h}{3}\left[f(x_0) + 4f(x_1) + f(x_2)\right]
\]

Como $h = \frac{b-a}{2}$, también puede escribirse como:

\[
\int_a^b f(x)\,dx \approx \frac{b-a}{6}
\left[f(a) + 4f\left(\frac{a+b}{2}\right) + f(b)\right]
\]

Esta es la \textbf{regla de Simpson 1/3}.



\subsection{Deducción del error de Simpson 1/3}

El error de Simpson proviene del error de interpolación cuadrática.

Sabemos que el error del polinomio interpolador de grado 2 es:

\[
f(x) - p_2(x) = \frac{f^{(3)}(\xi_x)}{3!}(x-x_0)(x-x_1)(x-x_2)
\]

El error de la integral es:

\[
E_S = \int_a^b [f(x) - p_2(x)] dx
\]

\[
E_S = \int_a^b \frac{f^{(3)}(\xi_x)}{6}
(x-x_0)(x-x_1)(x-x_2)\,dx
\]

Aplicando el teorema del valor medio para integrales:

\[
E_S = \frac{f^{(3)}(\xi)}{6}
\int_a^b (x-x_0)(x-x_1)(x-x_2)\,dx
\]

Al resolver esta integral (cálculo algebraico), se obtiene:

\[
E_S = -\frac{(b-a)^5}{2880} f^{(4)}(\xi)
\]

Por lo tanto, el error de Simpson es:

\[
E_S = -\frac{(b-a)^5}{180} f^{(4)}(\xi)
\]

\subsubsection*{Orden del método}

Si definimos $h = \frac{b-a}{2}$, entonces:

\[
E_S = O(h^4)
\]

Por lo tanto, la regla de Simpson es de orden 4.

%%%%%%%%%%%%%%%%%%%%%%%%%%%%%%%%
\subsection{Ejemplo Simpson 1/3}

\[
\int_0^1 x^2 dx
\]

\[
S = \frac{1}{6} [f(0) + 4f(0.5) + f(1)]
\]

\[
S = \frac{1}{6} [0 + 4(0.25) + 1] = \frac{2}{6} = \frac{1}{3}
\]

Simpson es exacto para polinomios de grado $\leq 3$.

\subsection{Simpson 3/8}

\[
\int_a^b f(x)\,dx \approx \frac{3h}{8} \left[f(x_0) + 3f(x_1) + 3f(x_2) + f(x_3)\right]
\]

Error:

\[
E = -\frac{(b-a)^5}{80} f^{(4)}(\xi)
\]

\newpage


\section{Regla Compuestas}

\section{Regla del Trapecio Compuesta}

\subsection{Deducción de la fórmula del trapecio compuesto}

Sea la integral:

\[
\int_a^b f(x)\,dx
\]

Dividimos el intervalo $[a,b]$ en $n$ subintervalos de igual tamaño:

\[
h = \frac{b-a}{n}
\]

Los puntos de la partición son:

\[
x_i = a + ih, \quad i = 0,1,2,\dots,n
\]

Aplicamos la regla del trapecio en cada subintervalo:

\[
\int_{x_i}^{x_{i+1}} f(x)\,dx \approx \frac{h}{2}\left[f(x_i) + f(x_{i+1})\right]
\]

Sumamos todas las aproximaciones:

\[
\int_a^b f(x)\,dx \approx \sum_{i=0}^{n-1} \frac{h}{2}\left[f(x_i) + f(x_{i+1})\right]
\]

Factorizando:

\[
= \frac{h}{2} \left[
f(x_0) + 2f(x_1) + 2f(x_2) + \dots + 2f(x_{n-1}) + f(x_n)
\right]
\]

Esta es la \textbf{regla del trapecio compuesta}.

\subsection{Deducción del error del trapecio compuesto}

Sabemos que el error del trapecio simple en un intervalo $[x_i, x_{i+1}]$ es:

\[
E_i = -\frac{h^3}{12} f''(\xi_i)
\]

donde $\xi_i \in (x_i, x_{i+1})$.

El error total es la suma de los errores en todos los subintervalos:

\[
E_T = \sum_{i=0}^{n-1} E_i
\]

\[
E_T = \sum_{i=0}^{n-1} \left(-\frac{h^3}{12} f''(\xi_i)\right)
\]

\[
E_T = -\frac{h^3}{12} \sum_{i=0}^{n-1} f''(\xi_i)
\]

Como $f''(x)$ es continua, podemos aplicar el teorema del valor medio y obtenemos:

\[
E_T = -\frac{h^3}{12} n f''(\xi)
\]

para algún $\xi \in (a,b)$.

Como $n = \frac{b-a}{h}$, entonces:

\[
E_T = -\frac{h^3}{12} \frac{b-a}{h} f''(\xi)
\]

\[
E_T = -\frac{(b-a)}{12} h^2 f''(\xi)
\]

Por lo tanto, el error del trapecio compuesto es:

\[
E_T = -\frac{(b-a)}{12} h^2 f''(\xi)
\]

\subsubsection*{Orden del método}

\[
E_T = O(h^2)
\]

Esto significa que si reducimos el tamaño de paso a la mitad, el error se reduce aproximadamente por un factor de 4.


\subsection{Ejemplo}

Aproximar la integral:

\[
\int_0^1 x^2 dx
\]

Sabemos que el valor exacto es:

\[
\int_0^1 x^2 dx = \frac{1}{3} = 0.3333
\]

Tomamos $n=4$ subintervalos:

\[
h = \frac{1-0}{4} = 0.25
\]

Los puntos son:

\[
x_0 = 0,\quad x_1 = 0.25,\quad x_2 = 0.5,\quad x_3 = 0.75,\quad x_4 = 1
\]

Evaluamos la función:

\[
f(x)=x^2
\]

\[
f(0)=0,\quad f(0.25)=0.0625,\quad f(0.5)=0.25,\quad f(0.75)=0.5625,\quad f(1)=1
\]

Aplicamos trapecio compuesto:

\[
\int_0^1 f(x)\,dx \approx \frac{h}{2}
\left[
f(x_0) + 2f(x_1) + 2f(x_2) + 2f(x_3) + f(x_4)
\right]
\]

\[
= \frac{0.25}{2}
\left[
0 + 2(0.0625) + 2(0.25) + 2(0.5625) + 1
\right]
\]

\[
= 0.125 (2.75)
\]

\[
= 0.34375
\]

Error:

\[
\text{Error} = |0.3333 - 0.34375| = 0.01045
\]


%%%%%%%%%%%%%%%%%%%%%%%%%%%

\subsection{Deducción de la fórmula de Simpson compuesta}

La regla de Simpson simple aproxima la integral en un intervalo usando una parábola.
Para mejorar la precisión, dividimos el intervalo $[a,b]$ en $n$ subintervalos iguales, con $n$ par.

Definimos el tamaño de paso:

\[
h = \frac{b-a}{n}, \quad n \text{ par}
\]

Los puntos son:

\[
x_i = a + ih, \quad i=0,1,2,\dots,n
\]

Aplicamos Simpson en cada par de subintervalos:

\[
\int_{x_0}^{x_2} f(x)\,dx \approx \frac{h}{3}[f(x_0) + 4f(x_1) + f(x_2)]
\]

\[
\int_{x_2}^{x_4} f(x)\,dx \approx \frac{h}{3}[f(x_2) + 4f(x_3) + f(x_4)]
\]

\[
\int_{x_4}^{x_6} f(x)\,dx \approx \frac{h}{3}[f(x_4) + 4f(x_5) + f(x_6)]
\]

Sumando todas las aproximaciones:

\[
\int_a^b f(x)\,dx \approx \frac{h}{3}
\left[
f(x_0)
+ 4f(x_1)
+ 2f(x_2)
+ 4f(x_3)
+ 2f(x_4)
+ \dots
+ 4f(x_{n-1})
+ f(x_n)
\right]
\]

Esta es la \textbf{regla de Simpson compuesta}.

\subsection{Error de Simpson compuesto}

Sabemos que el error de Simpson simple es:

\[
E = -\frac{(b-a)^5}{180} f^{(4)}(\xi)
\]

En Simpson compuesto, aplicamos Simpson en múltiples subintervalos, por lo que el error total es la suma de los errores en cada subintervalo.

Si el intervalo total se divide en $n$ subintervalos de tamaño $h$, el error total es:

\[
E_S = -\frac{b-a}{180} h^4 f^{(4)}(\xi)
\]

Por lo tanto, el error del método de Simpson compuesto es:

\[
E_S = O(h^4)
\]

Esto significa que si reducimos el tamaño de paso a la mitad, el error disminuye aproximadamente por un factor de $16$.


\subsection{Ejemplo}

Aproximar la integral:

\[
\int_0^1 x^4 dx
\]

Sabemos que el valor exacto es:

\[
\int_0^1 x^4 dx = \frac{1}{5} = 0.2
\]

Tomamos $n=4$ subintervalos:

\[
h = \frac{1-0}{4} = 0.25
\]

Los puntos son:

\[
x_0 = 0,\quad x_1 = 0.25,\quad x_2 = 0.5,\quad x_3 = 0.75,\quad x_4 = 1
\]

Evaluamos la función:

\[
f(x) = x^4
\]

\[
f(0)=0,\quad f(0.25)=0.0039,\quad f(0.5)=0.0625,\quad f(0.75)=0.3164,\quad f(1)=1
\]

Aplicamos Simpson compuesto:

\[
\int_0^1 f(x)dx \approx \frac{h}{3}
\left[
f(x_0) + 4f(x_1) + 2f(x_2) + 4f(x_3) + f(x_4)
\right]
\]

\[
= \frac{0.25}{3}
\left[
0 + 4(0.0039) + 2(0.0625) + 4(0.3164) + 1
\right]
\]

\[
= \frac{0.25}{3} (2.4062)
\]

\[
= 0.2005
\]

Error:

\[
\text{Error} = |0.2 - 0.2005| = 0.0005
\]


\newpage

\section{Cuadratura de Gauss}

\subsection{Idea del método}

A diferencia de las reglas de Newton-Cotes (trapecio, Simpson), la cuadratura de Gauss
no utiliza puntos equiespaciados. En su lugar, el método elige puntos y pesos óptimos
para maximizar la precisión.

La integral se aproxima por:

\[
\int_a^b f(x)\,dx \approx \sum_{i=1}^{n} w_i f(x_i)
\]

donde:
\begin{itemize}
\item $x_i$ son los puntos de Gauss
\item $w_i$ son los pesos
\end{itemize}

El método de Gauss con $n$ puntos es exacto para polinomios de grado $2n-1$.

\subsection{Deducción de la cuadratura de Gauss con 2 puntos}

Primero transformamos la integral al intervalo estándar $[-1,1]$:

\[
\int_a^b f(x)\,dx =
\frac{b-a}{2} \int_{-1}^{1} f\left(\frac{b-a}{2}t + \frac{a+b}{2}\right) dt
\]

La cuadratura de Gauss con 2 puntos tiene la forma:

\[
\int_{-1}^{1} f(t)\,dt \approx w_1 f(t_1) + w_2 f(t_2)
\]

Queremos que la fórmula sea exacta para polinomios hasta grado 3:

\[
\int_{-1}^{1} 1\,dt = 2
\]

\[
\int_{-1}^{1} t\,dt = 0
\]

\[
\int_{-1}^{1} t^2\,dt = \frac{2}{3}
\]

\[
\int_{-1}^{1} t^3\,dt = 0
\]

Entonces imponemos:

\[
w_1 + w_2 = 2
\]

\[
w_1 t_1 + w_2 t_2 = 0
\]

\[
w_1 t_1^2 + w_2 t_2^2 = \frac{2}{3}
\]

\[
w_1 t_1^3 + w_2 t_2^3 = 0
\]

Por simetría:

\[
t_1 = -t_2, \quad w_1 = w_2
\]

Entonces:

\[
2w_1 = 2 \Rightarrow w_1 = w_2 = 1
\]

\[
2w_1 t_1^2 = \frac{2}{3}
\]

\[
t_1^2 = \frac{1}{3}
\]

\[
t_1 = \pm \frac{1}{\sqrt{3}}
\]

Por lo tanto:

\[
\int_{-1}^{1} f(t)\,dt \approx
f\left(-\frac{1}{\sqrt{3}}\right)
+
f\left(\frac{1}{\sqrt{3}}\right)
\]

Finalmente, regresando al intervalo $[a,b]$:

\[
\int_a^b f(x)\,dx \approx
\frac{b-a}{2}
\left[
f\left(\frac{a+b}{2} - \frac{b-a}{2\sqrt{3}}\right)
+
f\left(\frac{a+b}{2} + \frac{b-a}{2\sqrt{3}}\right)
\right]
\]

Esta es la fórmula de Gauss con 2 puntos.

\subsection{Error de la cuadratura de Gauss}

La cuadratura de Gauss con $n$ puntos es exacta para polinomios de grado $2n-1$.
El error viene dado por:

\[
E_G = \frac{f^{(2n)}(\xi)}{(2n)!} C
\]

donde $C$ es una constante que depende del método.

Para Gauss con 2 puntos:

\[
E_G = -\frac{(b-a)^5}{90} f^{(4)}(\xi)
\]

Por lo tanto, el método de Gauss con 2 puntos tiene orden:

\[
E_G = O(h^4)
\]

Sin embargo, Gauss es más preciso que Simpson porque utiliza puntos óptimos.

\subsection{Ejemplo}

Aproximar la integral:

\[
\int_0^1 x^2 dx
\]

Usamos Gauss con 2 puntos.

Primero transformamos los puntos:

\[
x_1 = \frac{a+b}{2} - \frac{b-a}{2\sqrt{3}}
\]

\[
x_2 = \frac{a+b}{2} + \frac{b-a}{2\sqrt{3}}
\]

Como $a=0$, $b=1$:

\[
x_1 = \frac{1}{2} - \frac{1}{2\sqrt{3}} = 0.2113
\]

\[
x_2 = \frac{1}{2} + \frac{1}{2\sqrt{3}} = 0.7887
\]

Evaluamos:

\[
f(x_1) = (0.2113)^2 = 0.0447
\]

\[
f(x_2) = (0.7887)^2 = 0.6220
\]

Aplicamos Gauss:

\[
\int_0^1 x^2 dx \approx \frac{1}{2} [0.0447 + 0.6220]
\]

\[
= 0.3333
\]

Que coincide con el valor exacto:

\[
\int_0^1 x^2 dx = \frac{1}{3}
\]

La cuadratura de Gauss es exacta para polinomios de grado hasta 3.

\end{document}